\section{AutoResearch 方法的起源}

\subsection{Karpathy 的原始方法}

Andrej Karpathy（OpenAI 联合创始人、特斯拉前 AI 负责人、``vibe coding''一词的发明者）发布了一套叫做 \textbf{autoresearch} 的方法。

\begin{importantbox}{AutoResearch 核心思路}
不让你手动改进，而是让 AI agent 在循环里替你做：
\begin{enumerate}
    \item 尝试一个小改动
    \item 检查结果是否变好
    \item 变好了就留下，没变好就扔掉
    \item 然后再来一次，再来一次……
\end{enumerate}
\end{importantbox}

Karpathy 最初把这套方法用在机器学习代码上。但关键洞察在于：这套方法适用于\textbf{任何你能衡量并改进的东西}——包括你在 Claude 里搭建的 skills。

\subsection{作者的改造}

作者将 Karpathy 的方法改造成了一个可以在 Claude Code 和 Cowork 里运行的 skill。使用方式极其简单：

\begin{quote}
``run autoresearch on my landing page skill''
\end{quote}

说一句话，agent 就会接管整个优化过程。

\subsection{本章小结}

AutoResearch 源自 Karpathy，核心是``小改动 $\to$ 评估 $\to$ 保留或撤销''的自动循环。作者将其从机器学习场景推广到 Claude Skills 优化，封装成可直接调用的 skill。

